% Generated draft: requires mathematical and source review.
\documentclass[11pt]{article}
\usepackage[margin=25mm]{geometry}
\usepackage{fontspec}
\setmainfont{Latin Modern Roman}
\usepackage{amsmath,amssymb,mathtools,unicode-math}
\setmathfont{Latin Modern Math}
\usepackage{enumitem,needspace,xurl,hyperref}
\setlength{\parindent}{0pt}
\setlength{\parskip}{5pt}
\setlength{\emergencystretch}{3em}
\newcommand{\N}{\mathbb{N}}
\newcommand{\Z}{\mathbb{Z}}
\newcommand{\Q}{\mathbb{Q}}
\newcommand{\R}{\mathbb{R}}
\newcommand{\C}{\mathbb{C}}
\newcommand{\F}{\mathbb{F}}
\newcommand{\A}{\mathbb{A}}
\newcommand{\PP}{\mathbb{P}}
\newcommand{\E}{\mathbb{E}}
\newcommand{\eps}{\varepsilon}
\newcommand{\dd}{\,\mathrm{d}}
\begin{document}
\section*{UnsolvedMath ID: 30005641}
\textbf{Source problem code:} catalogue entry
\subsection{Definitions and mathematical statement}
The supplied record does not define the ambient object \(X\), the parameters \(\ell\) and \(d\), the group \(G_\ell\), the object \(V_\ell\), or the loci \(|X|_{V_\ell}\) and \(|X|_{\leq d}\). Accordingly, the following is a conditional statement in the source notation, requiring the missing definitions and hypotheses before it constitutes a fully specified conjecture.

For the intended \(X\), endowed with its Zariski topology, write
\[
S_{\ell,d}:=|X|_{V_\ell}\cap |X|_{\leq d}.
\]
Here both loci are retained as source-designated subsets of points of \(X\); their membership conditions are not supplied. A subset \(S\) is non-Zariski-dense in \(X\) if its Zariski closure \(\overline{S}^{\,\mathrm{Zar}}\) is a proper subset of \(X\). Write \(G_\ell^{\mathrm{ab}}\) for the abelianization intended by the source; the record does not specify whether this uses an abstract commutator subgroup or its closure in a group topology.

\textbf{Conjecture 1, as recorded.} Under the intended ambient hypotheses and for the intended admissible parameters, assume that \(G_\ell\) has finite abelianization. Then
\[
\#G_\ell^{\mathrm{ab}}<\infty
\quad\Longrightarrow\quad
\overline{|X|_{V_\ell}\cap |X|_{\leq d}}^{\,\mathrm{Zar}}\subsetneq X.
\]
The ambient hypotheses, parameter ranges, and quantifiers over \(X\), \(V_\ell\), \(\ell\), and \(d\) are absent from the supplied excerpt and must be recovered before this schema can be made precise. In particular, neither the meaning of the bound indexed by \(d\) nor the relationship between \(G_\ell\) and \(V_\ell\) is specified.

The supplied excerpt attributes the curve case to Cadoret--Tamagawa, with reference label [CT13], and describes the case \(\dim X\geq 2\) as widely open. It also reports that the strategy of [CT13] gives a heuristic for \(d=1\), referring to [Cad23, \S 4]. The subsequent explanation is truncated after a reference to a diophantine consideration. These are supplied source qualifications, not independently verified mathematical or literature assessments. An extracted footnote says that an unspecified object is only well defined up to conjugacy; its referent cannot be recovered from this record.
\subsection{Short English statement}
The recorded conjecture says that, if \(G_\ell\) has finite abelianization, the intersection of the two designated loci is not Zariski-dense in \(X\). The record omits the definitions of these loci and the ambient hypotheses, so the statement remains a schema requiring review. Supplied curation labels the problem partially solved, attributing the curve case to Cadoret--Tamagawa and describing dimensions at least two as widely open; these claims have not been independently verified.
\subsection{Sources}
\begin{itemize}
\item Supplied local UnsolvedMath registry record, permanent ID \(30005641\), exact source problem code \texttt{catalogue entry}, title ``Non-Density of Bounded-Rank Galois Loci,'' in the Oberwolfach Reports Open Problems collection. The record gives proposed year 2023, status \texttt{partially\_solved}, and statement status \texttt{corrected\_verified}. Its assertion that the statement was checked against the official report is supplied provenance, not an independent verification here.
\item Supplied report citation: \emph{MFO--RIMS Tandem Workshop: Arithmetic Homotopy and Galois Theory} (2024), Oberwolfach Reports, pp. 2377--2488, DOI \texttt{10.4171/owr/2023/42}. Supplied source URL: \url{https://doi.org/10.4171/owr/2023/42}. The extracted contribution is titled ``On the toric locus of \(\ell\)-adic local systems arising from geometry,'' Anna Cadoret, joint work with Jakob Stix. No separate proposer is supplied.
\item Supplied literature source: ``Cadoret--Stix report with the conjecture and known curve case,'' \url{https://doi.org/10.4171/owr/2023/42}. Literature checked at: \texttt{2026-08-22}. This date and the accompanying partially solved assessment are supplied dated curation, not independently verified. The background explicitly requests independent verification of time-sensitive status and cited claims.
\item Internal references preserved from the excerpt: [CT13], attributed in the supplied curation to Cadoret--Tamagawa for the curve case; and [Cad23, \S 4], cited for a heuristic when \(d=1\). Full bibliographic entries and separate links for these references are not supplied.
\end{itemize}
\end{document}
